\documentclass{article}
\usepackage[utf8]{inputenc}
\usepackage[T1]{fontenc}
\usepackage{geometry}
\geometry{
    a4paper,
    left=2.5cm,
    right=2.5cm,
    top=2.5cm,
    bottom=2.5cm
}
\usepackage{graphicx}
\usepackage{float}
\usepackage{hyperref}
\usepackage[english]{babel}
\usepackage{amsmath}
\usepackage{amssymb}
\usepackage{array}
\usepackage{booktabs}
\usepackage{caption}
\usepackage{subcaption}
\usepackage{enumitem}
\usepackage{fancyhdr}
\usepackage{titlesec}
\usepackage{url}
\usepackage{textcomp}
\usepackage{xcolor}
\usepackage{listings}
\usepackage{multirow}
\usepackage{longtable}
\usepackage{tabularx}
\usepackage{times}

% Set up page style
\pagestyle{fancy}
\fancyhf{}
\rfoot{\thepage}
\renewcommand{\headrulewidth}{0pt}

% Title formatting
\titleformat{\chapter}[display]{\normalfont\bfseries}{\chaptertitle\thechapter}{0pt}{\Huge}
\titlespacing*{\chapter}{0pt}{50pt}{40pt}

% Document info
\title{Template Document}
\author{Academic Support Agent}
\date{\today}

\begin{document}
\maketitle

\begin{figure}[h!]
\centering
\includegraphics[scale=0.18]{ufsj.jpg}
\end{figure}
%--------------------------------------------------------------
%--------------------------------------------------------------%--------------------------------------------------------------
%--------------------------------------------------------------
%	           LOGO DA UFSJ, CAPA E FOLHA DE ROSTO
%--------------------------------------------------------------
%--------------------------------------------------------------%--------------------------------------------------------------
%--------------------------------------------------------------
\begin{center}

{\fontsize{14}{0} \selectfont Departamento de Ci ncias Naturais\\
\vspace{0.15cm}
Programa de P s-Gradua o em F sica e Qu mica de Materiais}

\vspace{4cm}

{\fontsize{14}{0} \selectfont \textbf{Nome do Aluno}}

\vspace{5cm}

%       \vfill

{\fontsize{16}{0} \selectfont Titulo do projeto}\\

\vfill

%Department Name\\
%University Name\\
%Country\\
{\fontsize{12}{0} \selectfont S o Jo o del-Rei, M s de 20XX}\\
%{\fontsize{12}{0} \selectfont M s de 2020}
\end{center}
%--------------------------------------------------------------
%--------------------------------------------------------------%--------------------------------------------------------------
%--------------------------------------------------------------
\newpage
%--------------------------------------------------------------
%--------------------------------------------------------------%--------------------------------------------------------------
%--------------------------------------------------------------
\begin{center}

\vspace*{-1cm}
%\textbf{Thesis Title}
%\vspace{0.5cm}

{\fontsize{14}{0} \selectfont Universidade Federal de S o Jo o del-Rei\\
\vspace{0.3cm}
Departamento de Ci ncias Naturais\\
\vspace{0.3cm}
Programa de P s-Gradua o em F sica e Qu mica de Materiais}

\vspace{6cm}

%       \vfill

{\fontsize{16}{0} \selectfont Titulo do projeto}\\
\vspace{0.15cm}
% para mais escrita do t tulo \\
\vspace{0.15cm}
% para mais escrita do t tulo \\

\vspace{6cm}

\begin{flushright}
%\fbox									% Para preencher texto em forma de caixa
\begin{minipage}[b][4.5cm][c]{9.7cm}	% Configura o da caixa
Texto apresentado ao Programa de P s-gradua o em F sica e Qu mica de Materiais, da Universidade Federal de S o Jo o del-Rei, como parte do Exame de Qualifica o para obten o do t tulo de Mestre em F sica e Qu mica de Materiais.\\
\textbf{ rea de Concentra o:} F sica e Qu mica de Materiais\\
\textbf{Orientador:} \\
\textbf{Coorientador(a):} (Caso haja)
\end{minipage}
\end{flushright}



\vfill

%Department Name\\
%University Name\\
%Country\\

{\fontsize{12}{0} \selectfont S o Jo o del-Rei, M s de 20XX}\\
%{\fontsize{12}{0} \selectfont M s de 2020}
\end{center}
%--------------------------------------------------------------
%--------------------------------------------------------------%--------------------------------------------------------------
%--------------------------------------------------------------
%	                     -- SUM RIO
%--------------------------------------------------------------
%--------------------------------------------------------------%--------------------------------------------------------------
%--------------------------------------------------------------
% Content from contents would be included here

%--------------------------------------------------------------
%--------------------------------------------------------------%--------------------------------------------------------------
%--------------------------------------------------------------
%	                    -- INTRODU O
%	              editar no arquivo Tex : "intro"
%--------------------------------------------------------------
%--------------------------------------------------------------%--------------------------------------------------------------
%--------------------------------------------------------------
% Content from intro would be included here

%--------------------------------------------------------------
%--------------------------------------------------------------%--------------------------------------------------------------
%--------------------------------------------------------------
%	                    -- OBJETIVOS
%	              editar no arquivo Tex : "obje"
%--------------------------------------------------------------
%--------------------------------------------------------------%--------------------------------------------------------------
%--------------------------------------------------------------
% Content from obje would be included here

%--------------------------------------------------------------
%--------------------------------------------------------------%--------------------------------------------------------------
%--------------------------------------------------------------
%	                    -- METODOLOGIA
%	              editar no arquivo Tex : "meto"
%--------------------------------------------------------------
%--------------------------------------------------------------%--------------------------------------------------------------
%--------------------------------------------------------------
% Content from meto would be included here

%--------------------------------------------------------------
%--------------------------------------------------------------%--------------------------------------------------------------
%--------------------------------------------------------------
%	              -- RESULTADOS E DISCUSS O
%	              editar no arquivo Tex : "resul"
%--------------------------------------------------------------
%--------------------------------------------------------------%--------------------------------------------------------------
%--------------------------------------------------------------
% Content from resu would be included here

%--------------------------------------------------------------
%--------------------------------------------------------------%--------------------------------------------------------------
%--------------------------------------------------------------
%	              -- CONSIDERA ES FINAIS
%	              editar no arquivo Tex : "cons"
%--------------------------------------------------------------
%--------------------------------------------------------------%--------------------------------------------------------------
%--------------------------------------------------------------
% Content from cons would be included here

%--------------------------------------------------------------
%--------------------------------------------------------------%--------------------------------------------------------------
%--------------------------------------------------------------
%	                 -- PERSPECTIVAS
%	              editar no arquivo Tex : "pers"
%--------------------------------------------------------------
%--------------------------------------------------------------%--------------------------------------------------------------
%--------------------------------------------------------------
% Content from pers would be included here

%--------------------------------------------------------------
%--------------------------------------------------------------%--------------------------------------------------------------
%--------------------------------------------------------------
%	                  -- REFER NCIAS
%	    a gera o   autom tica com a utiliza o do BibTex
%--------------------------------------------------------------
%--------------------------------------------------------------%--------------------------------------------------------------
%--------------------------------------------------------------
%\renewcommand{\refname}{7 \hspace{2mm} Refer ncias}


%--------------------------------------------
\end{document}